\documentclass[10pt,a4paper]{article}
\usepackage[margin=19mm]{geometry}
\usepackage[T1]{fontenc}
\usepackage{lmodern}
\usepackage{microtype}
\usepackage{amsmath,amssymb,amsthm,mathtools}
\usepackage{graphicx}
\usepackage{xcolor}
\usepackage[hidelinks]{hyperref}

\definecolor{navy}{RGB}{21,55,94}
\definecolor{softblue}{RGB}{238,245,252}
\newtheorem{theorem}{Theorem}
\newtheorem{lemma}{Lemma}
\newtheorem{remark}{Remark}
\newcommand{\C}{\mathbb C}
\newcommand{\dd}{\,\mathrm d}
\newcommand{\dbar}{\bar\partial}
\newcommand{\ip}[2]{\langle #1,#2\rangle}

\title{\vspace{-1.4em}\textbf{A several-variable growing-weight
$\boldsymbol{\dbar}$ estimate}}
\author{Candidate full-result packet for arXiv:1311.2020}
\date{17 August 2026}

\begin{document}
\maketitle
\vspace{-1.4em}

\begin{center}
\fcolorbox{navy}{softblue}{\begin{minipage}{0.91\linewidth}
\textbf{Outcome.} Hedenmalm's Theorem 1.2 extends to every $\C^n$ with
the same factor $1/2$.  The Laplacian becomes the trace of the Levi form,
and the holomorphic-moment condition becomes annihilation against
anti-holomorphic-divergence-free coefficient vectors.  Positivity of the
Levi trace everywhere suffices; no uniform lower bound is required.
\end{minipage}}
\end{center}

\section*{1. The exact question and answer}

For $z=(z_1,\ldots,z_n)$ write
\[
 \phi_{j\bar j}=\partial_j\dbar_j\phi,
 \qquad L_\phi=\sum_{j=1}^n\phi_{j\bar j}.
\]
The source proves the following on $\C$: if $\Delta\phi>0$ and
$f\in L^2(e^{2\phi})$ annihilates every entire
$g\in L^2(e^{-2\phi})$, then $\dbar u=f$ and
\[
 \int |u|^2e^{2\phi}\Delta\phi
 \le \frac12\int |f|^2e^{2\phi}.
\]
Remark 1.3 asks for a several-variable generalization.

\begin{figure}[h]
\centering
\includegraphics[width=0.96\linewidth]{figures/open_problem_crop.png}
\caption{Theorem 1.2 and the several-variable question, source PDF page 2.}
\end{figure}

Define $\mathcal Z_\phi$ to be the coefficient vectors
$g=(g_1,\ldots,g_n)$ such that
\[
 \int_{\C^n}|g|^2e^{-2\phi}\dd V<\infty,
 \qquad \sum_{j=1}^n\dbar_jg_j=0
 \quad\hbox{in distributions}.
\]

\begin{theorem}[Growing weights on $\C^n$]\label{thm:main}
Let $\phi\in C^2(\C^n;\mathbb R)$ satisfy $L_\phi>0$ everywhere.  Let
$f=\sum_{j=1}^nf_j\dd\bar z_j$ obey
\[
 \int_{\C^n}|f|^2e^{2\phi}\dd V<\infty,
 \qquad
 \int_{\C^n}\sum_{j=1}^nf_jg_j\dd V=0
 \quad(g\in\mathcal Z_\phi).
 \tag{C}
\]
Then there is $u\in L^2_{\mathrm{loc}}(\C^n)$ with $\dbar u=f$ in
distributions and
\begin{equation}\label{eq:estimate}
 \int_{\C^n}|u|^2e^{2\phi}L_\phi\dd V
 \le \frac12\int_{\C^n}|f|^2e^{2\phi}\dd V.
\end{equation}
For $n=1$, this is exactly Theorem 1.2 of the source.
\end{theorem}

\paragraph{Proof intuition.}
Conjugate the equation by the growing weight.  The resulting operator $T$
has two decisive properties.  First, the difference of two squares in the
source's one-variable identity sums over the coordinates and controls a
scalar by $2L_\phi$.  Second, the kernel of $T^*$ becomes exactly
$\mathcal Z_\phi$ after undoing the weight and taking a conjugate.  Condition
(C) therefore puts the datum in the closure of the range of $T$; coercivity
turns a range-approximating sequence into an actual locally convergent
solution.

\section*{2. The conjugated identity}

On $C_c^\infty(\C^n)$ define the vector-valued operator
\[
 Tv=(T_1v,\ldots,T_nv),\qquad
 T_jv=\dbar_jv-v\dbar_j\phi
      =e^\phi\dbar_j(e^{-\phi}v).
\]

\begin{lemma}[Trace-Levi coercivity]\label{lem:coercive}
For every $v\in C_c^\infty(\C^n)$,
\begin{align}
 \|Tv\|_2^2
 &-\sum_{j=1}^n\|\partial_jv+v\partial_j\phi\|_2^2
   =2\int_{\C^n}|v|^2L_\phi\dd V,\label{eq:identity}\\
 2\int_{\C^n}|v|^2L_\phi\dd V&\le\|Tv\|_2^2.\label{eq:coercive}
\end{align}
\end{lemma}

\begin{proof}
For a fixed $j$, integration by parts in $z_j$ gives the source's identity
\[
 \|\dbar_jv-v\dbar_j\phi\|_2^2
 -\|\partial_jv+v\partial_j\phi\|_2^2
 =2\int_{\C^n}|v|^2\phi_{j\bar j}\dd V.
\]
The remaining variables are merely parameters.  Sum over $j$ to obtain
\eqref{eq:identity}; discard the nonnegative square to obtain
\eqref{eq:coercive}.
\end{proof}

\section*{3. Compatibility is the adjoint kernel}

Let $H_0=L^2(\C^n)$ and $H_1=L^2(\C^n;\C^n)$, both unweighted, and regard
$T_0=T|_{C_c^\infty}$ as a densely defined map $H_0\to H_1$.
Distributionally,
\begin{equation}\label{eq:adjoint}
 T_0^*k=-\sum_{j=1}^n(\partial_j+\partial_j\phi)k_j.
\end{equation}
If $k\in H_1$ and $g_j=e^\phi\overline{k_j}$, then
\[
 \int|g|^2e^{-2\phi}\dd V=\|k\|_2^2,
 \qquad
 \overline{T_0^*k}
 =-e^{-\phi}\sum_{j=1}^n\dbar_jg_j.
\]
Thus $k\mapsto g$ is a bijection from $\ker T_0^*$ onto
$\mathcal Z_\phi$.

Put $h=e^\phi(f_1,\ldots,f_n)\in H_1$.  With the standard complex Hilbert
pairing, for the corresponding $g$ we have
\[
 \ip{h}{k}_{H_1}
 =\int_{\C^n}\sum_{j=1}^nf_jg_j\dd V.
\]
Condition (C) is consequently equivalent to $h\perp\ker T_0^*$.  The
elementary Hilbert-space identity
\begin{equation}\label{eq:range}
 (\ker T_0^*)^\perp=\overline{\operatorname{ran}T_0}^{\,H_1}
\end{equation}
shows that there are $v_m\in C_c^\infty(\C^n)$ with $T_0v_m\to h$ in
$H_1$.

\section*{4. Closure produces the solution}

Apply \eqref{eq:coercive} to $v_m-v_\ell$.  Then
\[
 \int|v_m-v_\ell|^2L_\phi\dd V
 \le\frac12\|T_0v_m-T_0v_\ell\|_2^2,
\]
so $(v_m)$ converges to some $v$ in $L^2(L_\phi\dd V)$ and
\begin{equation}\label{eq:vbound}
 \int|v|^2L_\phi\dd V\le\frac12\|h\|_2^2.
\end{equation}
Because $L_\phi$ is continuous and strictly positive, it has a positive
minimum on each compact set.  Hence $v_m\to v$ also in $L^2_{\rm loc}$.
Passing to distributions in $T_0v_m\to h$ gives $Tv=h$ exactly.

Set $u=e^{-\phi}v$.  The intertwining formula defining $T$ yields
\[
 e^\phi\dbar u=Tv=h=e^\phi f,
\]
so $\dbar u=f$.  Finally, \eqref{eq:vbound} is precisely
\eqref{eq:estimate}.  This proves Theorem~\ref{thm:main}.

\section*{5. Why the condition must couple components}

A componentwise copy of the one-variable moment condition cannot work.
Indeed, choose $a\in C_c^\infty(\C^n)$ with
$\dbar_2\dbar_1a\not\equiv0$ and set
$f=(\dbar_1a,0,\ldots,0)$.  Each component annihilates all entire scalar
tests by integration by parts, yet $\dbar f\ne0$, so no scalar $u$ can have
$\dbar u=f$.  The divergence-free vectors in (C) are the exact coupled
replacement: they describe the full adjoint kernel and implicitly encode
the usual $\dbar$-compatibility.

When $n=1$, $\mathcal Z_\phi=A^2(\C,e^{-2\phi})$, because
$\dbar g=0$ distributionally means that $g$ is entire.  Also
$L_\phi=\Delta\phi$ in the source's normalized convention.  Thus both the
hypothesis and the constant reduce verbatim to Theorem 1.2.

\paragraph{Upgrade and novelty audit.}
Eight focused stages were completed: exact-signal and formulation audits;
rejection of componentwise moments; conjugated-operator reformulation;
trace-Levi identity; adjoint-kernel identification; closure construction;
and reduction, constant, and literature checks.  Bounded primary-source
searches through 17 August 2026 found no arXiv paper stating this exact
growing-weight result.  The Hilbert-space mechanism is elementary enough
that the result may be folklore; priority requires expert review.

\begin{thebibliography}{9}\small
\bibitem{Hed}
H.~Hedenmalm,
\emph{On H\"ormander's solution of the $\dbar$-equation},
arXiv:1311.2020 (2013).

\bibitem{Hor}
L.~H\"ormander,
\emph{An Introduction to Complex Analysis in Several Variables},
3rd ed., North-Holland, 1990.
\end{thebibliography}

\end{document}
